% Generated draft: requires mathematical and source review.
\documentclass[11pt]{article}
\usepackage[margin=25mm]{geometry}
\usepackage{fontspec}
\setmainfont{Latin Modern Roman}
\usepackage{amsmath,amssymb,mathtools,unicode-math}
\setmathfont{Latin Modern Math}
\usepackage{enumitem,needspace,xurl,hyperref}
\setlength{\parindent}{0pt}
\setlength{\parskip}{5pt}
\setlength{\emergencystretch}{3em}
\newcommand{\N}{\mathbb{N}}
\newcommand{\Z}{\mathbb{Z}}
\newcommand{\Q}{\mathbb{Q}}
\newcommand{\R}{\mathbb{R}}
\newcommand{\C}{\mathbb{C}}
\newcommand{\F}{\mathbb{F}}
\newcommand{\A}{\mathbb{A}}
\newcommand{\PP}{\mathbb{P}}
\newcommand{\E}{\mathbb{E}}
\newcommand{\eps}{\varepsilon}
\newcommand{\dd}{\,\mathrm{d}}
\begin{document}
\section*{UnsolvedMath ID: 701425}
\textbf{Source problem code:} catalogue entry
\subsection{Definitions and mathematical statement}
For a group \(G\), let \(G'=[G,G]\) be its commutator subgroup and let \(G''=[G',G']\). The group \(G\) is \emph{metabelian} if \(G''=\{1\}\), equivalently if \(G'\) is abelian. A finite presentation is an expression
\[
P=\langle x_1,\ldots,x_n\mid r_1,\ldots,r_m\rangle,
\]
where \(n,m\) are finite nonnegative integers and each \(r_j\) is a word in the generators and their inverses. It defines the group \(G_P=F_n/\langle\!\langle r_1,\ldots,r_m\rangle\!\rangle\), where \(F_n\) is the free group on the indicated generators and the denominator is the normal closure of the relators.

Fix a standard effective encoding of finite presentations by natural numbers. A set of codes is \emph{recursively enumerable} if an algorithm lists exactly its members, with repetitions allowed. The question, under the usual interpretation through finite presentations, is whether
\[
\mathcal M=\{\operatorname{code}(P): P\text{ is a finite presentation and }G_P''=\{1\}\}
\]
is recursively enumerable. Equivalently, can an algorithm enumerate finite presentations of metabelian groups so that every finitely presented metabelian group occurs up to isomorphism and no group outside this class occurs?

The source does not explicitly specify an encoding or an enumeration convention. This formulation uses the conventional effective presentation interpretation and does not require distinct isomorphism types to occur only once. The supplied registry labels the problem open while recording an explicitly unverified AI claim of an affirmative answer; that claim is not treated here as an established result.
\subsection{Short English statement}
Can an algorithm list finite presentations for all finitely presented metabelian groups, covering every such group up to isomorphism and producing only metabelian groups? Repetitions are allowed. The supplied open label and unverified affirmative solution claim are source qualifications requiring review.
\subsection{Sources}
\begin{itemize}
\item Local UnsolvedMath registry, supplied record with permanent ID \(701425\), exact source code \texttt{catalogue entry}, title \emph{Kourovka Notebook Problem 17.124}. The record attributes the problem to V. Shpilrain, proposed in 2010.
\item \emph{The Kourovka Notebook (Unsolved Problems in Group Theory)}, 17th issue (2010), Problem 17.124; also identified in the supplied record with the 21st issue edition (2026), \texttt{arXiv:1401.0300v48}. The supplied collection description names E. I. Khukhro and V. D. Mazurov as editors. These source identifications are supplied attribution and have not been independently checked.
\item The supplied discussion reports: ``[Unverified AI claim] Yes, it is'', attributed to Achyuth Jayadevan, \emph{Preprint}, 2026, \url{https://arxiv.org/abs/2609.10281}. No verification of this claim is supplied.
\item The supplied status review is dated \texttt{2026-10-16}, retains the canonical label \texttt{open}, and classifies the evidence as \texttt{unrefereed\_claim}. It reports a \(\#\) marker for an unverified AI solution claim in the October 2026 21st issue, \texttt{arXiv:1401.0300v48}. This assessment is supplied dated curation, not independently verified literature assessment.
\end{itemize}
\end{document}
